\section{Langages réguliers et analyse lexicale}
The chapter on analyse syntaxique ended by narrowing grammaires non
contextuelles until nothing much was left: les CFG où toutes les productions
sont de la forme
\[
  N \;\to\; \Sigma \times N \;\cup\; \epsilon ,
\]
that is, a terminal followed by a single non-terminal, or nothing at all. That
class has a name --- \emph{les langages réguliers}. This chapter is that class
seen from three sides at once: as expressions régulières, which a human
writes; as automates, which a machine runs; and as the front end of the
compilateur which uses both, the \textit{lexeur}. It closes on writing a lexeur
from scratch.

\subsection{Langages réguliers}
\begin{definition}[Diverses définitions]
  On peut définir les langages réguliers de diverses manières :
  \begin{itemize}
    \item via les \textbf{opérations rationnelles} (\textit{rational
      operations}) ;
    \item par les expressions régulières ;
    \item par les automates ;
    \item comme image inverse de \textbf{monoïdes finis} (\textit{finite
      monoids}) ;
    \item par la \textbf{logique} (\textit{logic}) ;
    \item etc.
  \end{itemize}
\end{definition}

\begin{remark}
  The list is not padding. A class of languages with one definition is a
  curiosity; a class with five unrelated-looking definitions that all pick out
  the same sets is a robust object, and each definition is good at a different
  job. The right-linear grammars above are a sixth entry, reached from the side
  of analyse syntaxique. Only two of them are developed here, and they are
  chosen for opposite reasons: expressions because they are what one writes,
  automates because they are what one runs.
\end{remark}

\subsection{Expressions régulières}
\subsubsection{Syntaxe}
\begin{definition}[Expressions régulières, inductivement]
  On définit inductivement les expressions régulières :
  \[
    \begin{array}{r@{\;}l@{\qquad}l}
      r \mathrel{\Coloneqq} & \emptyset      & \text{langage vide --- the empty language} \\
                            & \epsilon       & \text{mot vide --- the empty word} \\
                            & c              & c \text{ un caractère --- a character} \\
                            & [c_1 - c_2]    & \text{une classe de caractères --- a character class} \\
                            & r \mid r       & \text{union} \\
                            & r\,r           & \text{concaténation} \\
                            & r{*}           & \text{répétition (zéro ou plus)} \\
                            & r{+}           & \text{répétition (au moins une fois)}
    \end{array}
  \]
\end{definition}

Eight cases, of which the first four are atoms and the last four are
constructors. Note that this is a definition of \emph{syntax} only: at this
stage $r{*}$ is a tree node with one child, nothing more. What it denotes comes
next, and is a separate definition.

\begin{remark}
  $\emptyset$ and $\epsilon$ are not the same expression and not the same
  language. $\emptyset$ denotes the language with no words at all;
  $\epsilon$ denotes the language whose single word is the word of length zero.
  Concatenating anything with $\emptyset$ gives $\emptyset$; concatenating
  anything with $\epsilon$ changes nothing. The distinction matters as soon as
  one writes a lexeur rule that can match nothing.
\end{remark}

\subsubsection{Le langage associé}
\begin{definition}[Langage associé]
  On peut aussi définir le \textbf{langage associé} (\textit{associated
  language}) :
  \[
    \begin{array}{r@{\quad}l}
      \text{Expression } r & \text{Langage } L(r) \\[2pt]
      \emptyset    & \emptyset \\
      \epsilon     & \{\epsilon\} \\
      c            & \{c\} \\
      {}[c_1 - c_2]  & \{ c \mid c_1 \le c \le c_2 \} \\
      r_a \mid r_b & L(r_a) \union L(r_b) \\
      r_a r_b      & \{ w_a w_b \mid w_a \in L(r_a),\, w_b \in L(r_b) \} \\
      r{*}         & \{ w_1 \ldots w_k \mid k \ge 0 \wedge \forall i,\, w_i \in L(r) \} \\
      r{+}         & \{ w_1 \ldots w_k \mid k \ge 1 \wedge \forall i,\, w_i \in L(r) \}
    \end{array}
  \]
\end{definition}

The whole table is one structural recursion: the language of a compound
expression is built out of the languages of its parts, so $L$ is total and
unambiguous on the syntax above.

\begin{remark}
  The star and the plus differ in exactly one character of the definition,
  $k \ge 0$ against $k \ge 1$, so $r{+}$ is redundant: $L(r\,r{*}) = L(r{+})$.
  It is kept because it is the case one actually writes --- an integer literal
  is $[0-9]{+}$, never $[0-9][0-9]{*}$ --- and an implementation that omits
  the constructor simply builds it, writing \texttt{Concat(digit, Star(digit))}
  for \texttt{digits}, as in the examples below.
\end{remark}

\begin{remark}
  The character class $[c_1 - c_2]$ is defined by the \emph{order} on
  characters, $c_1 \le c \le c_2$: it is a range in the character encoding, not
  an arbitrary set. That is why \texttt{[a-z]} and \texttt{[A-Z]} have to be
  written as two classes joined by a union --- in ASCII they are not contiguous.
\end{remark}

\subsubsection{Des exemples}
The examples are given not as chaînes de caractères in concrete syntax but as
constructor calls: an expression is a value of an algebraic data type with one
constructor per line of the grammar.

\begin{lstlisting}[language=Python]
digit = CharRange('0','9') # one digit
digits = Concat(digit, Star(digit))
letter = Union(CharRange('a','z'),CharRange('A','Z'))
let_tok = Concat(Char('l'), Concat(Char('e'),Char('t')))
in_tok = Concat(CharRange('i','i'),CharRange('n','n'))
opp = Char('+')
opf = Char('*')
opm = Char('-')
ope = Char('=')
ident = Concat(letter,Star(Union(digit,letter)))
space = Union(Char(' '),Char('\t'))
\end{lstlisting}

These are the very definitions that the lexeur written at the end of this
chapter starts from. Two details of the encoding are worth
reading. \texttt{let\_tok} shows that concatenation is binary, so a three-letter
keyword is a nest of two \texttt{Concat}s; the surface syntax \texttt{let} hides
an association choice that the data type must make explicit. And
\texttt{in\_tok} uses \texttt{CharRange('i','i')} where \texttt{Char('i')} would
do --- a degenerate class of one character, harmless because $c \le c \le c$.

\begin{remark}
  The constructor set is \texttt{Char}, \texttt{CharRange}, \texttt{Union},
  \texttt{Concat}, \texttt{Star}: five of the eight cases of the grammar. There
  is no constructor for $\emptyset$, none for $\epsilon$ and none for $r{+}$, so
  an implementation of the lexeur only has to handle five cases --- with $r{+}$
  spelled \texttt{Concat(r, Star(r))} as in \texttt{digits}.
\end{remark}

\subsection{Les reconnaître}
Knowing what $L(r)$ \emph{is} does not say how to test membership. A
recogniser can work directly on the syntax tree, by mutual recursion over a
fixed word.

\begin{definition}[La fonction match]
  Étant donné un mot $w_0 \ldots w_{n-1}$ fixé on définit
  \[
    \operatorname{match}(R, i) \;=\; \{\, j \mid w_i \ldots w_{j-1} \in L(R) \,\}
  \]
  --- the set of positions at which a factor of $w$ starting at $i$ and matched
  by $R$ may end. On a alors :
  \[
    \begin{array}{r@{\;}c@{\;}l}
      \operatorname{match}(\epsilon, i)      &=& \{i\} \\[2pt]
      \operatorname{match}(c, i)             &=& \{ i+1 \mid \text{si } w_i = c \} \\[2pt]
      \operatorname{match}([c_1, c_2], i)    &=& \{ i+1 \mid \text{si } w_i \in [c_1, c_2] \} \\[2pt]
      \operatorname{match}(R_a \mid R_b, i)  &=& \operatorname{match}(R_a, i) \union \operatorname{match}(R_b, i) \\[2pt]
      \operatorname{match}(R_a / R_b, i)     &=& \bigcup_{j \in \operatorname{match}(R_a, i)} \operatorname{match}(R_b, j) \\[2pt]
      \operatorname{match}(R{*}, i)          &=& \text{point fixe}_X \Big( \{i\} \cup_{j \in X} \operatorname{match}(R, j) \Big)
    \end{array}
  \]
\end{definition}

\begin{notation}
  The fifth line is the concatenation case: it writes $R_a / R_b$ where the
  grammar writes $R_a R_b$, the slash standing in for the invisible
  concatenation operator. One may also write $\mathit{matches}(R, w, j)$ with
  the word passed explicitly; since $w$ is fixed throughout, it is the same
  function as $\operatorname{match}(R, j)$. Note also that the word is indexed
  here from $0$, as $w_0 \ldots w_{n-1}$, whereas the automate and the lexeur
  below index it from $1$, as $w_1 \ldots w_n$. Neither convention is wrong;
  only mixing them inside one implementation is.
\end{notation}

The shape of these equations is what makes them implementable. Every case
returns a \emph{set} of end positions, not a yes/no answer, because an
expression régulière can match several prefixes of the same suffix at once, and the appelant
--- the lexeur --- is going to want the largest of them. Union takes the union
of the two sets. Concatenation is a relational composition: run $R_a$ from $i$,
and from every place it can stop, run $R_b$.

\begin{example}[Several answers at once]
  Take $w = \mathtt{let}$ (so $n = 3$) and
  $R = \mathtt{ident} = \mathit{letter}\,(\mathit{digit} \mid \mathit{letter}){*}$.
  Then $\operatorname{match}(\mathit{letter}, 0) = \{1\}$, and starring the tail
  from position $1$ reaches $1$, $2$ and $3$, so
  \[
    \operatorname{match}(R, 0) = \{1, 2, 3\},
  \]
  because \texttt{l}, \texttt{le} and \texttt{let} are all identifiers. The
  recogniser answers ``it matches, and here is every place it could stop''; it
  is the lexeur that will pick $3$.
\end{example}

\begin{remark}
  The star case is stated as a point fixe in $X$ rather than as a recursive
  call, and that is not stylistic. The naive reading --- ``$\operatorname{match}(R{*},i)$
  is $\{i\}$ together with $\operatorname{match}(R{*}, j)$ for every
  $j \in \operatorname{match}(R,i)$'' --- fails to terminate whenever
  $\epsilon \in L(R)$, since then $i \in \operatorname{match}(R,i)$ and the
  recursion calls itself at the same position for ever. Reading it as the
  \emph{least} point fixe of a monotone map on subsets of
  $\{0, \ldots, n\}$ fixes this: one starts from $\{i\}$ and saturates, adding
  end positions until nothing new appears. The set is finite and the map only
  grows, so the iteration stops after at most $n+1$ rounds.
\end{remark}

\begin{remark}
  Working directly on the expression tree is correct but not cheap: on a word of
  length $n$ each position may be visited once per node of the tree, and the
  concatenation case may fan out over a whole set of intermediate positions.
  This is precisely the inefficiency that the automate view removes.
\end{remark}

\subsection{Automates}
\begin{definition}[Automate fini déterministe (\textit{deterministic finite automaton})]
  Un \textbf{automate} (\textit{automaton}) fini déterministe c'est :
  \begin{itemize}
    \item un ensemble d'\textbf{états} (\textit{states}) $Q$ ;
    \item un ensemble d'\textbf{états finaux} (\textit{final states})
      $F \subseteq Q$ ;
    \item un \textbf{état initial} (\textit{initial state}) $i \in Q$ ;
    \item une \textbf{fonction de transition} (\textit{transition function})
      $\delta : Q \times \Sigma \to Q$.
  \end{itemize}
  Un mot $w_1 \ldots w_n$ est \textbf{accepté} (\textit{accepted}) par
  l'automate si $q_n \in F$ avec $q_0 = i$ et $q_{i+1} = \delta(q_i, w_{i+1})$.
\end{definition}

Acceptance is thus a single left-to-right pass with one variable holding the
état: read a character, replace the état by
$\delta(\text{state}, \text{character}), $ and at the end ask whether it is an
état final. The cost is exactly one transition lookup per character, whatever
the expression behind the automate looked like.

\begin{remark}
  $\delta$ is a \emph{function} on all of $Q \times \Sigma$, so it is total:
  every état has an outgoing transition for every character of the alphabet,
  and the run never gets stuck --- it simply may end outside $F$. In a drawing
  one does not draw the transitions that lead nowhere useful; they all go to an
  implicit non-final \textbf{état puits} (\textit{sink state}) that loops on
  every character. Determinism is the other half of the definition: one target
  état per pair, hence one run per word, hence no backtracking.
\end{remark}

\begin{example}[The two automates behind the example lexeur]
  The rules \texttt{[0-9]+} and \texttt{[a-z]([a-z]|[0-9])*} of the worked
  example below become these two automates. Double circles are états finaux;
  the état puits is not drawn.
  \begin{figure}[h!]
    \centering
    \begin{tikzpicture}[shorten >=1pt, node distance=2.9cm, on grid, auto,
      every state/.style={draw, minimum size=7mm, inner sep=1pt, font=\small},
      font=\small]
      \node[state, initial, initial text={}] (a0) {$q_0$};
      \node[state, accepting, right=of a0] (a1) {$q_1$};
      \path[->]
        (a0) edge node {$[0\text{-}9]$} (a1)
        (a1) edge [loop above] node {$[0\text{-}9]$} ();
      \node[below=1.05cm of a0, anchor=north, font=\footnotesize] {$[0-9]{+}$};

      \node[state, initial, initial text={}, right=5.6cm of a0] (b0) {$p_0$};
      \node[state, accepting, right=of b0] (b1) {$p_1$};
      \path[->]
        (b0) edge node {$[a\text{-}z]$} (b1)
        (b1) edge [loop above] node {$[a\text{-}z] \mid [0\text{-}9]$} ();
      \node[below=1.05cm of b0, anchor=north, font=\footnotesize]
        {$[a-z]([a-z] \mid [0-9]){*}$};
    \end{tikzpicture}
  \end{figure}
  Both have two useful états, and both spend one transition per character.
  Compare with running $\operatorname{match}$ over the syntax tree of the same
  expressions: the automate has already done that work, once, at construction
  time.
\end{example}

\subsubsection{Comparaison}
\begin{center}
  \begin{tabularx}{\textwidth}{@{}lX@{}}
    \toprule
    \textbf{Automates} & facile à manipuler pour l'ordinateur (union,
      intersection, test d'appartenance, etc.) ; difficile à comprendre pour un
      humain \\
    \midrule
    \textbf{Expressions régulières} & facile à comprendre pour un humain ; moins
      efficace algorithmiquement \\
    \bottomrule
  \end{tabularx}
\end{center}

\begin{conclusion}
  \emph{Ce que l'on veut : écrire des expressions régulières puis les compiler
  en automates !} One writes the notation that is readable and compiles it into
  the representation that is fast --- which is the same division of labour the
  whole module is about, one level down.
\end{conclusion}

\subsection{Analyse lexicale : le lexeur}
The \textit{lexeur} --- the analyseur lexical, the first pass of the
compilateur --- is the first real consumer of all of this.

\begin{definition}[Règles et actions]
  Une manière de faire des analyseurs lexicaux consiste à utiliser une liste
  d'expressions régulières $e_1 \ldots e_k$ et une liste d'actions
  $a_1 \ldots a_k$. Le lexeur va découper le mot en \textbf{lexèmes}
  $t_1 \ldots t_\ell$ et pour le lexème $t_j$ reconnu par $e_i$ il va faire
  l'action $a_i(t_j)$.
\end{definition}

\begin{figure}[h!]
  \centering
  \begin{tikzpicture}[font=\small, node distance=0.6cm]
    \node[draw, rounded corners=3pt, minimum height=1.5cm, text width=3.1cm,
          align=center] (rules)
      {list of pairs\\ $(e_i, a_i)_{i=1..k}$\\ \footnotesize expression, action};
    \node[draw, rounded corners=3pt, minimum height=1.5cm, text width=2.4cm,
          align=center, right=1.6cm of rules] (lex) {\textbf{lexeur}};
    \node[draw, rounded corners=3pt, minimum height=1.5cm, text width=3.1cm,
          align=center, right=1.6cm of lex] (out)
      {lexèmes $t_1 \ldots t_\ell$\\ \footnotesize with $a_i(t_j)$ performed};
    \node[above=0.75cm of lex, font=\footnotesize] (src) {the word $w$ (the source text)};
    \path[->, >=latex]
      (rules) edge (lex)
      (lex) edge (out)
      (src) edge (lex);
  \end{tikzpicture}
\end{figure}

\subsubsection{Un exemple}
One text is run through one rule table, in order. Texte :
\begin{center}
  \texttt{let y\ \ \ = 42 in zzz+1\ \ \ \ 2+my\_var}
\end{center}
and the rules, \emph{in this order}, are

\begin{center}
  \begin{tabular}{@{}ll@{}}
    \toprule
    Expression & Action($x$) \\
    \midrule
    $[0-9]{+}$                        & Afficher \texttt{INT}($x$) \\
    \texttt{let}                      & Afficher \texttt{LET} \\
    \texttt{in}                       & Afficher \texttt{IN} \\
    \texttt{=}                        & Afficher \texttt{EQ} \\
    \texttt{'\ '}                     & Ignorer \\
    \texttt{+}                        & Afficher \texttt{PLUS} \\
    $[a-z]([a-z] \mid [0-9]){*}$      & Afficher \texttt{VAR}($x$) \\
    \bottomrule
  \end{tabular}
\end{center}

\noindent
Affichage :
\[
  \begin{array}{l}
    \texttt{LET},\;
    \texttt{VAR("y")},\;
    \texttt{EQ},\;
    \texttt{INT("42")},\;
    \texttt{IN},\;
    \texttt{VAR("zzz")},\;
    \texttt{PLUS}, \\[2pt]
    \texttt{INT("1")},\;
    \texttt{INT("2")},\;
    \texttt{PLUS},\;
    \texttt{VAR("my")},\;
    \textbf{ERREUR}.
  \end{array}
\]

\begin{figure}[h!]
  \centering
  \begin{tikzpicture}[font=\ttfamily\small, node distance=0pt]
    \tikzset{
      tok/.style={draw, rounded corners=2pt, minimum height=0.62cm,
                  inner xsep=3pt, anchor=west},
      ign/.style={draw, densely dotted, rounded corners=2pt,
                  minimum height=0.62cm, inner xsep=3pt, anchor=west},
      lab/.style={font=\sffamily\scriptsize, rotate=90, anchor=east, yshift=-3pt},
    }
    \node[tok] (t1) at (0,0) {let};
    \node[ign, right=1pt of t1] (s1) {\textvisiblespace};
    \node[tok, right=1pt of s1] (t2) {y};
    \node[ign, right=1pt of t2] (s2) {\textvisiblespace\textvisiblespace\textvisiblespace};
    \node[tok, right=1pt of s2] (t3) {=};
    \node[ign, right=1pt of t3] (s3) {\textvisiblespace};
    \node[tok, right=1pt of s3] (t4) {42};
    \node[ign, right=1pt of t4] (s4) {\textvisiblespace};
    \node[tok, right=1pt of s4] (t5) {in};
    \node[ign, right=1pt of t5] (s5) {\textvisiblespace};
    \node[tok, right=1pt of s5] (t6) {zzz};
    \node[tok, right=1pt of t6] (t7) {+};
    \node[tok, right=1pt of t7] (t8) {1};
    \node[ign, right=1pt of t8] (s6) {\textvisiblespace\textvisiblespace\textvisiblespace\textvisiblespace};
    \node[tok, right=1pt of s6] (t9) {2};
    \node[tok, right=1pt of t9] (t10) {+};
    \node[tok, right=1pt of t10] (t11) {my};
    \node[tok, right=1pt of t11, draw=red, thick] (t12) {\_var};

    \node[lab] at (t1.south) {LET};
    \node[lab] at (t2.south) {VAR};
    \node[lab] at (t3.south) {EQ};
    \node[lab] at (t4.south) {INT};
    \node[lab] at (t5.south) {IN};
    \node[lab] at (t6.south) {VAR};
    \node[lab] at (t7.south) {PLUS};
    \node[lab] at (t8.south) {INT};
    \node[lab] at (t9.south) {INT};
    \node[lab] at (t10.south) {PLUS};
    \node[lab] at (t11.south) {VAR};
    \node[lab, text=red] at (t12.south) {ERREUR};
  \end{tikzpicture}
\end{figure}

Dotted boxes are the matches whose action is \emph{Ignorer}; they consume input
and emit nothing. The run stops at the underscore: no rule of the table matches
a prefix beginning with \verb|_|, so \texttt{my\_var} comes out as
\texttt{VAR("my")} followed by an error rather than as one identifier.

\subsubsection{L'algorithme, formellement}
\begin{definition}[La boucle du lexeur]
  Formellement, l'analyseur procède sur le mot $w = w_1 \ldots w_n$ de la façon
  suivante :
  \begin{itemize}
    \item s'il a fini de lire le mot ($n = 0$) il s'arrête ;
    \item sinon il regarde le \textbf{plus grand préfixe} $w_1 \ldots w_k$ du
      mot qui est reconnu par une des expressions $e_i$ ;
    \item si aucune expression ne reconnaît de préfixe, on déclenche une
      erreur ;
    \item parmi les expressions $e_i$ qui reconnaissent le préfixe
      $w_1 \ldots w_k$, il choisit celle avec le $i$ le \textbf{plus petit} ;
    \item il effectue $a_i(w_1 \ldots w_k)$ ;
    \item et continue sur le mot $w_{k+1} \ldots w_n$.
  \end{itemize}
\end{definition}

Two tie-breaks, applied in this order and only in this order. First the length:
the longest prefix matched by \emph{any} rule wins, whichever rule that is ---
le \textbf{plus grand préfixe} (\textit{maximal munch}). Only then, among the
rules that match that same longest prefix, the earliest in the list wins ---
la \textbf{priorité de la règle} (\textit{rule priority}). Everything a lexeur
specification does is arranged by these two.

\begin{example}[Priorité de la règle : why \texttt{let} is not a variable]
  At the start of the text, the prefix \texttt{let} is matched by rule $2$
  (\texttt{let}) and by rule $7$ ($[a-z]([a-z] \mid [0-9]){*}$), both of length
  $3$, and no rule matches anything longer --- the next character is a space.
  Length does not separate them; the priorité de la règle does, and rule $2$
  wins, so the lexème is \texttt{LET} and not \texttt{VAR("let")}. A keyword is
  simply a rule placed above the identifier rule.
\end{example}

\begin{example}[Plus grand préfixe : why \texttt{lettre} is not a keyword]
  Reverse the situation and feed the lexeur the text \texttt{lettre}. Rule $2$
  matches the prefix \texttt{let}, of length $3$; rule $7$ matches
  \texttt{lettre}, of length $6$. The first tie-break is the length, and $6 > 3$,
  so the priorité de la règle is never consulted: the lexème is
  \texttt{VAR("lettre")}. This is the whole reason the length test comes first
  --- with the priorité de la règle first, no identifier could ever begin with a
  keyword.
\end{example}

\begin{figure}[h!]
  \centering
  \begin{tikzpicture}[shorten >=1pt, node distance=2.35cm, on grid, auto,
    every state/.style={draw, minimum size=7mm, inner sep=1pt, font=\small},
    font=\small]
    \node[state, initial, initial text={}] (s) {$s$};
    \node[state, accepting, right=of s] (A) {$A$};
    \node[state, accepting, right=of A] (B) {$B$};
    \node[state, accepting, right=of B] (C) {$C$};
    \node[state, accepting, below=2.4cm of B] (I) {$I$};
    \path[->]
      (s) edge node {$\mathtt{l}$} (A)
      (A) edge node {$\mathtt{e}$} (B)
      (B) edge node {$\mathtt{t}$} (C)
      (s) edge [bend right=20] node[swap, pos=0.42] {other $[a\text{-}z]$} (I)
      (A) edge node[swap, pos=0.35] {other} (I)
      (B) edge node[pos=0.4] {other} (I)
      (C) edge [bend left=20] node[pos=0.42] {$[a\text{-}z]\mid[0\text{-}9]$} (I)
      (I) edge [loop below] node {$[a\text{-}z]\mid[0\text{-}9]$} ();
  \end{tikzpicture}
\end{figure}

The picture is the two rules \texttt{let} and $[a-z]([a-z]\mid[0-9]){*}$ merged
into one automate, which is what a lexeur generator actually builds; the edges
marked ``other'' carry every letter or digit except the one that continues the
keyword. The états $A$, $B$, $I$ are final for the identifier rule only, so
reaching them means \texttt{VAR}; $C$ is final for both rules at once, and it is
there that the priorité de la règle decides in favour of \texttt{LET}. The
scanner does not test the rules one after another --- it runs this single
automate, remembers the last état final it passed and which rules were
satisfied there, and on getting stuck falls back to that position. That is the
plus grand préfixe and the priorité de la règle, compiled.

\begin{remark}
  The rule \texttt{'\ '} of the table is a single quoted space --- one character,
  not one or more. Since its action is \emph{Ignorer}, a run of several spaces
  simply goes round the loop several times, which is why the four-space run in
  \texttt{zzz+1\ \ \ \ 2} costs four iterations and produces nothing (and the
  three-space gap before the \texttt{=} costs three). The
  \texttt{ocamllex} version below writes \texttt{[' ' '\textbackslash t'
  '\textbackslash n']+} instead, folding the repetition into the expression and
  also accepting tabs and newlines.
\end{remark}

\begin{remark}
  The error case is the only place the lexeur can fail, and it fails on a
  \emph{character}, not on a structure: it reports that no rule can start here.
  Everything about how the lexèmes fit together --- that \texttt{let} must be
  followed by a variable, that \texttt{+} needs two operands --- is beyond the
  lexeur's reach and is the parser's problem.
\end{remark}

\subsubsection{ocamllex}
Rather than write the loop by hand, one declares the rule table in a
\texttt{.mll} file and lets a generator emit the automate and the scanner.

\begin{lstlisting}[language=Caml]
{ open Lexing }

rule token = parse
  | ['0'-'9']+ as s { `INT (int_of_string s) }
  | "let"     { `LET }
  | "in"      { `IN }
  | '+'       { `PLUS }
  | '='       { `EQ }
  | [' ' '\t' '\n']+ { token lexbuf }
  | ['a'-'z']['a'-'z''0'-'9']* as id { `VAR id }
  | eof       { `EOF }
\end{lstlisting}

À mettre dans \texttt{lexer.mll} puis compiler avec \texttt{ocamllex lexer.mll}
--- \emph{utile pour votre compilateur !}\\

The correspondence with the table of the worked example is exact: each
\texttt{|} line is a pair $(e_i, a_i)$, the expression on the left in the
concrete syntax of the tool, the action on the right as an OCaml expression
between braces. Three points are specific to the tool.

\begin{itemize}
  \item \verb|as s| binds the matched text --- what the definition calls
    $w_1 \ldots w_k$, the argument $x$ of the action --- to a variable, so
    \texttt{INT} and \texttt{VAR} can carry it. Rules whose action ignores the
    text, such as \texttt{"let"}, need no binding.
  \item The whitespace rule's action is \texttt{token lexbuf}: it calls the
    scanner again on the rest of the input and returns whatever that returns.
    A recursive call \emph{is} how one writes \emph{Ignorer} when every rule has
    to produce a lexème.
  \item \texttt{eof} is an extra pattern the tool provides for end of input,
    turning the loop's termination condition into an ordinary rule; here it
    yields \texttt{`EOF}, a lexème the parser can match on.
\end{itemize}

\begin{remark}
  The results are backtick-prefixed --- \texttt{`INT}, \texttt{`LET} --- so they
  are OCaml \emph{polymorphic variants}: they need no type declared in advance,
  which is convenient for a type of lexèmes that grows as the compilateur does.
  The order of the lines is still load-bearing: \texttt{"let"} sits above the
  identifier rule for exactly the reason of the priorité de la règle example
  above.
\end{remark}

\subsubsection{Pourquoi un lexeur ?}
\begin{conclusion}
  Tout ce que le lexeur fait peut souvent être fait dans la grammaire mais ça :
  \begin{itemize}
    \item permet de décomposer le problème ;
    \item permet d'utiliser des expressions régulières ;
    \item simplifie beaucoup l'écriture et la lecture des grammaires ;
    \item diminue la complexité des grammaires.
  \end{itemize}
\end{conclusion}

The first clause is the honest one: a grammaire non contextuelle can generate
identifiers and integer literals character by character, so the lexeur is not a
necessity. It is a separation of concerns that pays. A grammar written over
lexèmes rather than characters has no rules for whitespace anywhere, does not
have to encode the plus grand préfixe as grammar rules, and --- since the
langages réguliers are a strict subclass of the languages of grammaires non
contextuelles --- hands the cheap part of the work to the cheap machine.

\subsection{Écrire un lexeur}
\begin{problem}[Objectif]
  Vous devez écrire du code qui implémente un petit lexeur, c'est-à-dire qui :
  \begin{itemize}
    \item récupère une liste de paires (expression régulière, nom) ;
    \item une chaîne de caractères ;
    \item et renvoie le résultat de l'analyse lexicale.
  \end{itemize}
\end{problem}

\noindent
\emph{Travailler avec méthode} : vous pouvez écrire les fonctions dans l'ordre
en testant chacune ; pensez à l'algorithme et s'il est correct ; réfléchissez à
ce qui est nécessaire au niveau du code.

\subsubsection{L'interface du lexeur}
Two extension files, \texttt{parser.ml} and \texttt{parser.py} --- the
continuation \emph{pour aller plus loin}, described in the last remark of this
chapter --- are written against the finished lexeur, so they pin down its
interface exactly. The rule list is the table of the worked example reordered
and slightly widened --- a rule for \texttt{*} is added, \texttt{INT} moves
below \texttt{ID}, \texttt{ident} also accepts capitals and \texttt{space} also
accepts a tab --- with a name in place of an action:

\begin{lstlisting}[language=Python]
rules = [
    (let_tok, "LET"), (in_tok, "IN"), (ope, "EQ"), (opp, "PLUS"),
    (opf, "FOIS"), (ident, "ID"), (digits, "INT"),
    (space, "SPACE")    # as in the worked example
]
\end{lstlisting}

\noindent
and it is consumed by a single function:

\begin{lstlisting}[language=Python]
list_tokens = lexer(rules, input_string)
rev_list_of_tokens = [("EOF","")] + list(reversed(
    [e for e in list_tokens if e[0] != "SPACE"]))
\end{lstlisting}

\begin{lstlisting}[language=Caml]
let token = ref ((lexer rules input_string
                  |> List.filter (fun (x,_) -> x != `SPACE))
                 @ [(`EOF,"")])
\end{lstlisting}

So \texttt{lexer} takes the rule list and the input chaîne de caractères and
returns a \emph{list of pairs} (name, matched text), in order of appearance ---
the OCaml version using polymorphic variants \texttt{`LET}, \texttt{`ID} as
names where the Python version uses chaînes de caractères. Three consequences
for the implementation.

\begin{itemize}
  \item The action of a rule is not a callback: it is just the name attached to
    the lexème, and the appelant decides what to do. \emph{Ignorer} is therefore
    not implemented in the lexeur at all --- \texttt{SPACE} lexèmes are produced
    like any other and the appelant filters them out, which is exactly what both
    snippets do before the analyse syntaxique.
  \item The matched text must be returned alongside the name, since the parser
    needs \texttt{"42"} to build the integer and the identifier's spelling to
    look it up in its environment (\texttt{Hashtbl.create 17} in OCaml,
    \texttt{dict()} in Python).
  \item The position of a rule in \texttt{rules} is the index $i$ of its
    priorité de la règle, so the list must be traversed in order: \texttt{LET}
    and \texttt{IN} precede \texttt{ID}, and an implementation that returns the
    first rule matching
    \emph{any} prefix, instead of the first rule matching the \emph{longest}
    prefix, will lex \texttt{lettre} as \texttt{LET} followed by
    \texttt{ID("tre")}.
\end{itemize}

\begin{remark}
  The natural shape of the code follows the two definitions of this chapter, in
  order: a \texttt{match} over the expression tree in the sense of the
  recognition equations, returning the set of end positions from a given start;
  then the loop of the lexeur, which at each position asks every rule for its
  end positions, keeps the largest $k$ over all rules, breaks ties by the
  smallest $i$, emits the pair and restarts at $k$. The error case is the one
  where every rule returns an empty set.
\end{remark}

\begin{remark}
  Une extension naturelle consiste à utiliser votre lexeur pour faire de
  l'analyse syntaxique. Vous pouvez utiliser les fichiers \texttt{parser.ml} ou
  \texttt{parser.py} à compléter pour avoir un parseur d'un langage contenant
  des entiers, des expressions \texttt{+}, \texttt{*} et des variables
  introduites par la syntaxe : \texttt{let \ldots = \ldots in \ldots}. The
  parser evaluates the expression to a number and fails if it is not parsable
  or uses an undefined variable. The skeleton they give --- a mutable list of
  lexèmes with \texttt{peek\_type} and \texttt{eat}, a hash table for the
  environment --- is a recursive-descent parser of exactly the kind the grammar
  chapter builds, which is why the grammar it is aimed at is the one that
  chapter asks to be made LL(1).
\end{remark}
